% AI Usage Declaration — TMA4268 Compulsory Exercise 2
% Include this file from report.tex with % AI Usage Declaration — TMA4268 Compulsory Exercise 2
% Include this file from report.tex with % AI Usage Declaration — TMA4268 Compulsory Exercise 2
% Include this file from report.tex with % AI Usage Declaration — TMA4268 Compulsory Exercise 2
% Include this file from report.tex with \input{ai-usage}

\section*{Declaration of AI Usage}

This project made use of AI coding assistants, specifically Claude (Anthropic),
during the development process.
The scope and boundaries of that usage are described below.

\subsection*{Work performed by the human authors}

All intellectual and analytical work in this project was carried out by the
human authors.
This includes, but is not limited to:

\begin{itemize}
  \item \textbf{Problem formulation.}
    Choosing the UFC dataset, framing the task as binary classification of the
    Red-corner winner, and deciding to benchmark against bookmaker-implied
    probabilities.

  \item \textbf{Exploratory data analysis.}
    Investigating the class imbalance (58.1\% Red wins), identifying the
    Red-corner assignment convention, examining missingness patterns,
    analysing the correlation structure among fighter statistics,
    and interpreting the violin-plot distributions.

  \item \textbf{Feature engineering decisions.}
    Designing the \emph{difference features} representation
    ($\Delta x = x_{\text{Red}} - x_{\text{Blue}}$), reasoning about its
    symmetry properties, deciding which contextual columns to retain
    (gender, title bout, number of rounds), and choosing to exclude
    odds-derived columns from the feature set to prevent leakage.

  \item \textbf{Model selection.}
    Choosing the three model families (logistic regression, Random Forest,
    XGBoost) and the Elastic Net variant with nested cross-validation.
    Selecting hyperparameter grids, deciding on the regularisation strategy,
    and setting tree-depth constraints (\texttt{min\_samples\_leaf},
    \texttt{min\_child\_weight}).

  \item \textbf{Evaluation design.}
    Designing the chronological 80/20 train-test split to prevent time
    leakage, selecting evaluation metrics (log loss, Brier score, ROC AUC,
    accuracy), and framing the bookmaker as the primary benchmark.

  \item \textbf{Interpretation of results.}
    Interpreting Elastic Net coefficients (e.g.\ the negative age coefficient
    and its implications), reading Random Forest feature importances,
    analysing confusion matrices, discussing the bias--variance trade-off
    across model families, and drawing conclusions about the bookmaker gap.

  \item \textbf{Report writing.}
    Drafting, structuring, and revising the full \LaTeX{} report, including
    the abstract, methodology descriptions, results discussion, and
    limitations section.
\end{itemize}

In short, all decisions that require understanding of statistical learning
--- choosing features, selecting and configuring models, designing the
evaluation protocol, and interpreting the output --- were made by the authors.

\subsection*{Work assisted by AI}

The AI assistant was used primarily as a \emph{coding tool} for implementation
tasks that do not involve statistical reasoning.
Typical prompts included instructions such as:

\begin{itemize}
  \item Plotting and formatting: ``change the colour of this graph'',
    ``create violin plots on \texttt{[X, Y, Z]} features'',
    ``add axis labels and a legend''.
  \item Pipeline assembly: ``make a pipeline with min-max scaling and median
    imputation'', ``wrap this in a \texttt{ColumnTransformer}''.
  \item Code scaffolding: ``generate a confusion matrix plot from these
    predictions'', ``save the metrics as JSON''.
  \item Repository organisation: setting up the \texttt{Makefile},
    configuring directory paths, managing imports, and structuring the
    Python package layout.
\end{itemize}

These are tasks where the human author specified \emph{what} to produce and
the AI translated that specification into syntactically correct code.
The AI did not choose which plots to make, which models to train, which
features to include, or how to interpret any result.

\subsection*{Summary}

The division of labour can be characterised as follows: the authors were
responsible for the \emph{statistical learning} content of the project
(analysis, modelling decisions, interpretation), while the AI assistant
served as a productivity tool for the \emph{software engineering} aspects
(code syntax, plotting boilerplate, file organisation).
No part of the written report was generated by AI.


\section*{Declaration of AI Usage}

This project made use of AI coding assistants, specifically Claude (Anthropic),
during the development process.
The scope and boundaries of that usage are described below.

\subsection*{Work performed by the human authors}

All intellectual and analytical work in this project was carried out by the
human authors.
This includes, but is not limited to:

\begin{itemize}
  \item \textbf{Problem formulation.}
    Choosing the UFC dataset, framing the task as binary classification of the
    Red-corner winner, and deciding to benchmark against bookmaker-implied
    probabilities.

  \item \textbf{Exploratory data analysis.}
    Investigating the class imbalance (58.1\% Red wins), identifying the
    Red-corner assignment convention, examining missingness patterns,
    analysing the correlation structure among fighter statistics,
    and interpreting the violin-plot distributions.

  \item \textbf{Feature engineering decisions.}
    Designing the \emph{difference features} representation
    ($\Delta x = x_{\text{Red}} - x_{\text{Blue}}$), reasoning about its
    symmetry properties, deciding which contextual columns to retain
    (gender, title bout, number of rounds), and choosing to exclude
    odds-derived columns from the feature set to prevent leakage.

  \item \textbf{Model selection.}
    Choosing the three model families (logistic regression, Random Forest,
    XGBoost) and the Elastic Net variant with nested cross-validation.
    Selecting hyperparameter grids, deciding on the regularisation strategy,
    and setting tree-depth constraints (\texttt{min\_samples\_leaf},
    \texttt{min\_child\_weight}).

  \item \textbf{Evaluation design.}
    Designing the chronological 80/20 train-test split to prevent time
    leakage, selecting evaluation metrics (log loss, Brier score, ROC AUC,
    accuracy), and framing the bookmaker as the primary benchmark.

  \item \textbf{Interpretation of results.}
    Interpreting Elastic Net coefficients (e.g.\ the negative age coefficient
    and its implications), reading Random Forest feature importances,
    analysing confusion matrices, discussing the bias--variance trade-off
    across model families, and drawing conclusions about the bookmaker gap.

  \item \textbf{Report writing.}
    Drafting, structuring, and revising the full \LaTeX{} report, including
    the abstract, methodology descriptions, results discussion, and
    limitations section.
\end{itemize}

In short, all decisions that require understanding of statistical learning
--- choosing features, selecting and configuring models, designing the
evaluation protocol, and interpreting the output --- were made by the authors.

\subsection*{Work assisted by AI}

The AI assistant was used primarily as a \emph{coding tool} for implementation
tasks that do not involve statistical reasoning.
Typical prompts included instructions such as:

\begin{itemize}
  \item Plotting and formatting: ``change the colour of this graph'',
    ``create violin plots on \texttt{[X, Y, Z]} features'',
    ``add axis labels and a legend''.
  \item Pipeline assembly: ``make a pipeline with min-max scaling and median
    imputation'', ``wrap this in a \texttt{ColumnTransformer}''.
  \item Code scaffolding: ``generate a confusion matrix plot from these
    predictions'', ``save the metrics as JSON''.
  \item Repository organisation: setting up the \texttt{Makefile},
    configuring directory paths, managing imports, and structuring the
    Python package layout.
\end{itemize}

These are tasks where the human author specified \emph{what} to produce and
the AI translated that specification into syntactically correct code.
The AI did not choose which plots to make, which models to train, which
features to include, or how to interpret any result.

\subsection*{Summary}

The division of labour can be characterised as follows: the authors were
responsible for the \emph{statistical learning} content of the project
(analysis, modelling decisions, interpretation), while the AI assistant
served as a productivity tool for the \emph{software engineering} aspects
(code syntax, plotting boilerplate, file organisation).
No part of the written report was generated by AI.


\section*{Declaration of AI Usage}

This project made use of AI coding assistants, specifically Claude (Anthropic),
during the development process.
The scope and boundaries of that usage are described below.

\subsection*{Work performed by the human authors}

All intellectual and analytical work in this project was carried out by the
human authors.
This includes, but is not limited to:

\begin{itemize}
  \item \textbf{Problem formulation.}
    Choosing the UFC dataset, framing the task as binary classification of the
    Red-corner winner, and deciding to benchmark against bookmaker-implied
    probabilities.

  \item \textbf{Exploratory data analysis.}
    Investigating the class imbalance (58.1\% Red wins), identifying the
    Red-corner assignment convention, examining missingness patterns,
    analysing the correlation structure among fighter statistics,
    and interpreting the violin-plot distributions.

  \item \textbf{Feature engineering decisions.}
    Designing the \emph{difference features} representation
    ($\Delta x = x_{\text{Red}} - x_{\text{Blue}}$), reasoning about its
    symmetry properties, deciding which contextual columns to retain
    (gender, title bout, number of rounds), and choosing to exclude
    odds-derived columns from the feature set to prevent leakage.

  \item \textbf{Model selection.}
    Choosing the three model families (logistic regression, Random Forest,
    XGBoost) and the Elastic Net variant with nested cross-validation.
    Selecting hyperparameter grids, deciding on the regularisation strategy,
    and setting tree-depth constraints (\texttt{min\_samples\_leaf},
    \texttt{min\_child\_weight}).

  \item \textbf{Evaluation design.}
    Designing the chronological 80/20 train-test split to prevent time
    leakage, selecting evaluation metrics (log loss, Brier score, ROC AUC,
    accuracy), and framing the bookmaker as the primary benchmark.

  \item \textbf{Interpretation of results.}
    Interpreting Elastic Net coefficients (e.g.\ the negative age coefficient
    and its implications), reading Random Forest feature importances,
    analysing confusion matrices, discussing the bias--variance trade-off
    across model families, and drawing conclusions about the bookmaker gap.

  \item \textbf{Report writing.}
    Drafting, structuring, and revising the full \LaTeX{} report, including
    the abstract, methodology descriptions, results discussion, and
    limitations section.
\end{itemize}

In short, all decisions that require understanding of statistical learning
--- choosing features, selecting and configuring models, designing the
evaluation protocol, and interpreting the output --- were made by the authors.

\subsection*{Work assisted by AI}

The AI assistant was used primarily as a \emph{coding tool} for implementation
tasks that do not involve statistical reasoning.
Typical prompts included instructions such as:

\begin{itemize}
  \item Plotting and formatting: ``change the colour of this graph'',
    ``create violin plots on \texttt{[X, Y, Z]} features'',
    ``add axis labels and a legend''.
  \item Pipeline assembly: ``make a pipeline with min-max scaling and median
    imputation'', ``wrap this in a \texttt{ColumnTransformer}''.
  \item Code scaffolding: ``generate a confusion matrix plot from these
    predictions'', ``save the metrics as JSON''.
  \item Repository organisation: setting up the \texttt{Makefile},
    configuring directory paths, managing imports, and structuring the
    Python package layout.
\end{itemize}

These are tasks where the human author specified \emph{what} to produce and
the AI translated that specification into syntactically correct code.
The AI did not choose which plots to make, which models to train, which
features to include, or how to interpret any result.

\subsection*{Summary}

The division of labour can be characterised as follows: the authors were
responsible for the \emph{statistical learning} content of the project
(analysis, modelling decisions, interpretation), while the AI assistant
served as a productivity tool for the \emph{software engineering} aspects
(code syntax, plotting boilerplate, file organisation).
No part of the written report was generated by AI.


\section*{Declaration of AI Usage}

This project made use of AI coding assistants, specifically Claude (Anthropic),
during the development process.
The scope and boundaries of that usage are described below.

\subsection*{Work performed by the human authors}

All intellectual and analytical work in this project was carried out by the
human authors.
This includes, but is not limited to:

\begin{itemize}
  \item \textbf{Problem formulation.}
    Choosing the UFC dataset, framing the task as binary classification of the
    Red-corner winner, and deciding to benchmark against bookmaker-implied
    probabilities.

  \item \textbf{Exploratory data analysis.}
    Investigating the class imbalance (58.1\% Red wins), identifying the
    Red-corner assignment convention, examining missingness patterns,
    analysing the correlation structure among fighter statistics,
    and interpreting the violin-plot distributions.

  \item \textbf{Feature engineering decisions.}
    Designing the \emph{difference features} representation
    ($\Delta x = x_{\text{Red}} - x_{\text{Blue}}$), reasoning about its
    symmetry properties, deciding which contextual columns to retain
    (gender, title bout, number of rounds), and choosing to exclude
    odds-derived columns from the feature set to prevent leakage.

  \item \textbf{Model selection.}
    Choosing the three model families (logistic regression, Random Forest,
    XGBoost) and the Elastic Net variant with nested cross-validation.
    Selecting hyperparameter grids, deciding on the regularisation strategy,
    and setting tree-depth constraints (\texttt{min\_samples\_leaf},
    \texttt{min\_child\_weight}).

  \item \textbf{Evaluation design.}
    Designing the chronological 80/20 train-test split to prevent time
    leakage, selecting evaluation metrics (log loss, Brier score, ROC AUC,
    accuracy), and framing the bookmaker as the primary benchmark.

  \item \textbf{Interpretation of results.}
    Interpreting Elastic Net coefficients (e.g.\ the negative age coefficient
    and its implications), reading Random Forest feature importances,
    analysing confusion matrices, discussing the bias--variance trade-off
    across model families, and drawing conclusions about the bookmaker gap.

  \item \textbf{Report writing.}
    Drafting, structuring, and revising the full \LaTeX{} report, including
    the abstract, methodology descriptions, results discussion, and
    limitations section.
\end{itemize}

In short, all decisions that require understanding of statistical learning
--- choosing features, selecting and configuring models, designing the
evaluation protocol, and interpreting the output --- were made by the authors.

\subsection*{Work assisted by AI}

The AI assistant was used primarily as a \emph{coding tool} for implementation
tasks that do not involve statistical reasoning.
Typical prompts included instructions such as:

\begin{itemize}
  \item Plotting and formatting: ``change the colour of this graph'',
    ``create violin plots on \texttt{[X, Y, Z]} features'',
    ``add axis labels and a legend''.
  \item Pipeline assembly: ``make a pipeline with min-max scaling and median
    imputation'', ``wrap this in a \texttt{ColumnTransformer}''.
  \item Code scaffolding: ``generate a confusion matrix plot from these
    predictions'', ``save the metrics as JSON''.
  \item Repository organisation: setting up the \texttt{Makefile},
    configuring directory paths, managing imports, and structuring the
    Python package layout.
\end{itemize}

These are tasks where the human author specified \emph{what} to produce and
the AI translated that specification into syntactically correct code.
The AI did not choose which plots to make, which models to train, which
features to include, or how to interpret any result.

\subsection*{Summary}

The division of labour can be characterised as follows: the authors were
responsible for the \emph{statistical learning} content of the project
(analysis, modelling decisions, interpretation), while the AI assistant
served as a productivity tool for the \emph{software engineering} aspects
(code syntax, plotting boilerplate, file organisation).
No part of the written report was generated by AI.
